%%%%%%%%%%%%%%%%%%%%%%preface.tex%%%%%%%%%%%%%%%%%%%%%%%%%%%%%%%%%%%%%%%%%
% sample preface
%
% Use this file as a template for your own input.
%
%%%%%%%%%%%%%%%%%%%%%%%% Springer %%%%%%%%%%%%%%%%%%%%%%%%%%

\preface

%% Please write your preface here
There is a long tradition regarding mathematical modeling of biological phenomena, which dates back to the very beginning of modern science. However, there has been a revival in the last decade and a half,  in which much effort has been made to elucidate the dynamic behavior of intracellular processes. This has been due in part to the development of impressive novel experimental techniques that permit precise measurements at the single cell, and even at the single molecule levels. But also thanks to the availability of low-cost and highly efficient computational power, which permits to numerically explore the dynamics of quite complex systems. One particular problem that has generated big interest is the stochasticity of gene expression. 

Over the last years it has become clear that the dynamics of most biological phenomena can be studied via the techniques of either nonlinear dynamics or stochastic processes. In either case, the biological system is usually visualized as a set of interdependent chemical reactions and the model equations are derived out of this picture. Deterministic, nonlinear dynamic models rely on chemical kinetics, while stochastic models are developed from the chemical master equation. Recent publications have demonstrated that deterministic models are nothing but an average description of the behavior of unicellular stochastic models. In that sense, the most detailed modeling approach is that of stochastic processes. However, both the deterministic and the stochastic approaches are complementary. The vast amount of available techniques to analytically explore the behavior of deterministic, nonlinear dynamical models is almost completely inexistent for their stochastic counterparts. On the other hand, the only way to investigate biochemical noise is via stochastic processes.

It is my impression that the excitement of developing a new science (that has has been termed systems biology) has made people unaware of a large amount of related discoveries performed in the late 19th and early 20th centuries by people like Planck, Nernst, de Groot, Prigogine etc. which we have classified as physical chemistry and irreversible thermodynamics. Moreover, although not as extensive a dismissal, there is no general acknowledgment on the part of systems biologists of the common mathematical grounds shared by biochemical and electrophysiological models; the later ones are usually classified in the realm of biophysics. One of the objectives of the present book is to show that the approaches of deterministic nonlinear dynamics, stochastic chemical master equation, and irreversible thermodynamical chemistry are all complementary, and that their proper combination allows a deeper understanding of the dynamic behavior of a large variety of biological systems. In particular, we tackle in this book gene expression and ion transport across membranes.

There are many great books on nonlinear dynamics, stochastic processes, irreversible thermodynamics, physical chemistry, and biochemistry. Nonetheless, to the best of my knowledge, there is none that brings all of these theories together, in an introductory but formal and comprehensive manner, for people interested in modeling biological phenomena. The present book is aimed at filling, at least in part, this gap. In order to achieve this goal I decided to follow a hands-on constructivist approach. The theory is developed stepwise, starting from the simplest concepts, and building upon them to derive, one step at a time, a more general framework. But instead of first developing the theory and latter studying its possible applications to biological systems, the examples are introduced right away. On the one hand, each theoretical development is motivated by specific biological examples. But also, every new mathematical derivation is immediately applied to one or more biological systems.

The target audience of this book are mainly last year undergraduate and graduate students with a solid mathematical background (phycisists, mathematicians, and engineers), as well as with basic notions of biochemistry and cellular biology, who are interested in learning the previously described techniques to model biological phenomena. The book can also be useful to students with a biological background, who are interested in mathematical modeling, and have a working knowledge of calculus, differential equations, and basic notions of probability theory.


Like all long term processes, the writing of the present book involved numerous people and institutions. In particular, I am in debt with my working place, the \emph{Centro de Investigaci\'on y de Estudios Avanzados del Instituto Polit\'ecnico Nacional}, for granting me a sabbatical leave in which I found the necessary time to write.
 I also thank the Department of Physiology of McGill University and the International Centre for Theoretical Physics for their hospitality and for providing me with the most adequate environment to carry out this project. Finally, of all the people who were involved, one way or another, in the process of writing the present book, I wish to emphasize Michael C. Mackey, Emanuel Salazar-Cabazos, Rom\'an U. Zapui\'en-Campos, and Luis U. Aguilera-de-Lira. I am deeply grateful with all the persons who contributed to the present book, but the support and advice of the above-mentioned people were so important that, in strict sense, made this book possible. In spite of several reviews it is possible that some mistakes are found along the book. I, and no one else, am the only responsible for all of them.

\vspace{\baselineskip}
\begin{flushright}\noindent
{\it Mois\'es Santill\'an\hfill } December, 2013\\
\end{flushright}


